\documentclass[11pt]{article}
\usepackage[margin=1in]{geometry}
\usepackage{amsmath,amssymb}
\usepackage{booktabs}
\usepackage{xcolor}
\newcommand{\todo}[1]{{\color{red}[TODO: #1]}}
\title{\textbf{LeWAM}: A Latent World-Action Model\\ \large When does latent world-model planning help for manipulation?}
\author{Minghao Fu \quad \ldots \quad Randall Balestriero}
\date{Draft skeleton --- \today}

\begin{document}
\maketitle

\begin{abstract}
We study \emph{when} a latent world model can plan for contact-rich manipulation. We introduce \textbf{LeWAM}, a latent world-action model that augments a JEPA latent world model with a jointly trained, history-conditioned action pathway, yielding a single model that both predicts latent dynamics and proposes actions. In contrast to the dominant world-action models---whether video-generative (DreamZero) or latent (Being-H0.7)---which learn to \emph{avoid} test-time search and amortize planning into one forward pass, LeWAM deliberately retains short-horizon planning and asks \emph{when it helps}. Our central result is an analysis of when latent world-model planning helps: short-horizon CEM against a goal-image latent cost succeeds on goal-reaching tasks where the goal encodes a large, visible, data-supported state change (OGBench-Cube, Push-T, Reacher), but collapses to near-zero on contact-rich robomimic (Lift/Can/Square), where planning fails for two distinct reasons we separate with probes: on Lift the success variable (grasp) is visually subdominant in a single frame, while on Can/Square the latent does see grasp but the task is multi-stage and a flat goal-image cost cannot decompose it. We show that a history / past-action-conditioned behavior-cloning pathway recovers this grasp state and far outperforms online planning on contact. Through a multi-training-seed study we further find that, at convergence, the action-head objective (detached vs.\ grounded decoder, raw-action vs.\ embedding-target, anti-collapse regularization) does not move success rate beyond training-run noise---a cautionary negative result. LeWAM's contribution is therefore the per-family planning analysis and the history-conditioned action pathway, not a state-of-the-art success rate.
\end{abstract}

\section{Introduction}
\textbf{From behavior cloning to world-action models.} \todo{1 para: VLA/BC map obs$\to$action with no explicit dynamics; the field is moving to world-action models (WAMs) that jointly model world + action, giving a forward model to plan/verify with. Two flavors: video-generative (DreamZero~\cite{dreamzero}) and latent (the LeWM~\cite{lewm} lineage). LeWAM is the latent WAM.}

\textbf{The question.} Latent world models plan by sampling-based search (CEM/MPC) against a latent goal cost. This works on some tasks and fails on others; \emph{when, and why?} \todo{frame the per-family question as the paper's spine.}

\textbf{Contributions.}
\begin{itemize}
  \item \textbf{A per-family analysis of latent WM-planning.} We show CEM planning succeeds on goal-reaching/geometry tasks and collapses to $\approx 0$ on contact-rich robomimic, and we identify the mechanism: the goal-image latent cost rewards the task only when the goal encodes a large, visible, data-supported state change (\S\ref{sec:perfamily}).
  \item \textbf{A history-conditioned action pathway} that far outperforms online planning on all contact tasks; on Lift it provably recovers grasp state that a single frame exposes only weakly (probe AUC $0.82$ vs $0.99$ from history), while on multi-stage Can/Square it supplies the sequential execution a flat planner lacks (\S\ref{sec:history}).
  \item \textbf{An honest multi-training-seed account}: at convergence the action-head objective does not change success rate beyond training-run noise; we make no SOTA-SR claim and instead position on mechanism and analysis (\S\ref{sec:negative}).
  \item \textbf{The LeWAM reframe}: a JEPA world model and its policy head as a single latent world-action model, the lightweight latent counterpart to video-generative WAMs.
\end{itemize}

% Related Work for LeWAM (Latent World-Action Model).
% Self-contained \section fragment; \input into the paper.
% Citations use readable keys; the key -> arXiv ID + title map is in the
% trailing comment block. Post-Jan-2026 IDs are flagged "verify".

\section{Related Work}
\label{sec:related}

\paragraph{World-action models: video-generative versus latent.}
A \emph{world-action model} (WAM) jointly models environment dynamics and the action that drives them, $p(\text{future},\,\text{action})$, so that a single learned object both predicts what happens next and supplies the control to make it happen. This frames the WAM against two neighbors. Unlike a vision-language-action policy~\cite{openvla,rt2}, which maps observations directly to actions with no explicit forward model, a WAM carries dynamics and can reason about counterfactual futures. Unlike a pure latent or 3D world model~\cite{dinowm,vjepa2}, which predicts futures but defers action selection to a separate inverse model or to test-time search, a WAM internalizes the action pathway. The family splits on the medium of the future. \emph{Video-generative} WAMs, exemplified by DreamZero~\cite{dreamzero}, model $p(\text{video},\,\text{action})$ with video diffusion, inherit physics priors from internet-scale video, and learn inverse dynamics from their own predicted rollouts so that no test-time optimization is required. \emph{Latent} WAMs, the LeWM~\cite{lewm} lineage we extend, predict in a self-supervised joint-embedding latent rather than in pixels, keeping the model light and the dynamics abstract. Being-H0.7~\cite{beingh07} is a recent large-scale latent WAM, organizing a future-aware latent world-action workspace from $200$k hours of egocentric video; like DreamZero it is search-free, discarding its training-only posterior branch and performing no rollout at inference. A common thread thus emerges: the dominant WAMs, whether video-generative (DreamZero) or latent (Being-H0.7), \emph{learn to avoid test-time search}, amortizing planning into a single forward pass. Fast-WAM~\cite{fastwam} makes the trend explicit, asking whether world-action models need test-time future imagination at all and answering \emph{no}---uniformly---finding that video co-training rather than test-time rollout drives control. LeWAM is also a latent WAM---a JEPA predictor with a jointly trained action pathway---but it revisits this question for latent models with \emph{explicit} search and answers it \emph{per task family} rather than uniformly: it deliberately retains a short-horizon search loop and asks \emph{when does latent world-model planning actually help, and when should one fall back to a history-conditioned policy?} Where Fast-WAM removes test-time imagination wholesale, LeWAM shows planning still pays off on goal-reaching/geometry tasks and collapses on contact-rich manipulation, and supplies the history-conditioned pathway that carries the latter case. Its contribution is therefore not to be a latent WAM (Being-H0.7 already is one at scale) but to characterize, per task family, where planning over a latent world model pays off and where it collapses---and to supply the history-conditioned pathway that carries the latter case.

\paragraph{JEPA and latent world models for control.}
A growing line of work learns a world model in a self-supervised joint-embedding (JEPA) latent and plans within it, dispensing with pixel reconstruction. LeWM~\cite{lewm}, the latent on which we build, trains a decoder-free predictor for pixel-latent consistency and plans by sampling-based optimization (CEM) against a latent goal-image cost; DINO-WM~\cite{dinowm} plans the same way over frozen DINO features with explicitly ``no pre-learned inverse models,'' and V-JEPA\,2-AC~\cite{vjepa2} extends a video JEPA to a real Franka arm, again with pure CEM at action time. While these systems convincingly show that a JEPA latent supports planning across navigation, deformables, and continuous control, none carries a learned action prior: the shared pattern is \emph{pretrained features plus sampling-based CEM/MPC with no action model}, which leaves every candidate rollout at the mercy of a dynamics model never trained to model contact. The field supplies its own evidence that this is brittle, V-JEPA\,2-AC's grasp success on a physical arm sitting well below its reaching performance, exactly the regime where rollouts must traverse contact dynamics the latent does not capture. A complementary line improves the latent's \emph{geometry} for planning rather than its action model: temporal straightening~\cite{tempstraight} regularizes the JEPA representation so that latent trajectories are locally straight and Euclidean cost approximates geodesic distance, sharpening the goal-image planning landscape and raising goal-reaching success. This is orthogonal to what we find: straightening conditions a geometry in which the task-relevant variable is already present, whereas our robomimic failures are of a different kind that no metric-geometry regularizer addresses---on Lift the success variable (grasp) is \emph{under-exposed} in the single frame, and on Can/Square it is present yet the task is \emph{multi-stage} so a flat goal-image cost cannot decompose it---both recovered instead by our history-conditioned pathway executing the demonstrated sequence. In contrast, LeWAM keeps this JEPA substrate but attaches a jointly trained action pathway and studies, per task family, where the substrate's planner does and does not work.

\paragraph{Learned policy as a planning prior.}
A second tradition warm-starts search with a learned policy, the ``thinking fast and slow'' division of labor introduced by expert iteration~\cite{exit}. TD-MPC2~\cite{tdmpc2} and its predecessor~\cite{tdmpc} learn a Markovian policy $\pi(a\,|\,z_t)$ on a decoder-free latent and use it to seed MPPI, and the MuZero family~\cite{muzero} plays the same role for discrete-action MCTS. Recent instances follow suit: policy-guided world-model planning for language-conditioned navigation~\cite{pijepa} seeds search from an instantaneous policy, and the JEPA-native Newt~\cite{newt} pairs a behavior-cloned prior with annealed-constraint planning across a massively multitask, padded-and-masked action space. These priors successfully turn an expensive search into a guided one, but across the cluster the prior is \emph{instantaneous and raw-action}: it conditions on the current latent state and proposes a single next action. The sole precedent that conditions on a past action, MBOP~\cite{mbop} with $a_t=f_b(s_t,a_{t-1})$, is feed-forward in one previous action, non-JEPA, and raw-action; DINO-WM and V-JEPA\,2-AC condition their \emph{dynamics} on action history yet carry no prior at all. In contrast, LeWAM's action pathway is \emph{autoregressive over past actions} $a_{<t}$ on a JEPA latent rather than Markovian, which is what lets it recover the contact state (``am I holding the object'') that a single-frame latent cannot expose. Newt and LeWAM both predict actions on JEPA latents; LeWAM is distinguished by the history and past-action conditioning and by the planning analysis that motivates it.

\paragraph{Amortized and search-free latent action.}
A third position replaces search outright by directly predicting actions in latent space. The most relevant instance is GC-IDM~\cite{gcidm}, a goal-conditioned inverse-dynamics model trained on the \emph{same LeWM latent}, which maps (current latent, goal latent, remaining horizon) to the next action and matches or exceeds CEM across navigation, manipulation, and control at roughly two orders of magnitude lower per-decision cost; Diffuser~\cite{diffuser} and Decision Diffuser~\cite{decisiondiffuser} likewise generate trajectories in one shot from return or skill guidance, and DPPO~\cite{dppo} makes the policy itself the planner. These methods compellingly show that amortization can match search at a fraction of the cost, but in making amortization the destination they forfeit any test-time correction when a single forward pass is wrong. In contrast, LeWAM treats amortization as a \emph{warm start to be refined}, with the action pathway proposing and short-horizon search correcting it where a single inverse-dynamics pass is brittle. GC-IDM, sharing our base, is the must-beat amortization baseline against which the warm-start-then-refine wedge has to earn its cost.

\paragraph{Hierarchical, diffusion-execution, and manipulation planners.}
On contact-rich manipulation specifically, the dominant recipe is to let a world model plan subgoals and a learned policy execute them. WorldDP~\cite{worlddp} diagnoses our setting directly, noting that visual world-model MPC is ``limited to single-stage tasks such as reaching or grasping'' and struggles with multi-stage ones, and its fix is hierarchical: a high-level world model that optimizes feasible subgoals at runtime with a low-level diffusion policy reaching each, over object-centric slot representations. LDP~\cite{ldp} is the closest published latent planner on our benchmark, reporting robomimic image Lift/Can/Square with a diffusion forecaster and diffusion inverse dynamics rather than a JEPA latent plus CEM, while FF-JEPA~\cite{ffjepa}, Diffusion Policy~\cite{diffusionpolicy}, and VQ-BeT~\cite{vqbet} round out the execution-side toolkit of expressive, multimodal action models. These planners convincingly solve multi-stage manipulation, but each assumes the failure of flat latent planning rather than mapping it. In contrast, LeWAM proposes no new hierarchy; its contribution on this axis is a per-family analysis of \emph{when flat latent planning fails}. Single-frame goal-image CEM succeeds when the goal encodes a large, visible, data-supported state change (e.g.\ an OGBench cube carried to a new pose) and collapses to near-zero on robomimic, where the success variable is visually subdominant and the demonstrations are narrow. This both explains why the hierarchical split helps and locates the failure precisely, on a contact-rich arena (robomimic Lift/Can/Square) that the world-model planning literature barely reports.

\paragraph{Representation versus planning.}
A sharper challenge questions the planning premise itself. MR.Q~\cite{mrq} argues that predictive \emph{representation} learning, not model-based planning, drives scalable multitask continuous control, exhibiting a model-free actor-critic with auxiliary predictive objectives that beats a world-model planner with no planning at all, and traces that gain by ablation to representation quality itself: stripping the predictive objective collapses the latent's effective rank and idles much of the critic, a shortfall of representational capacity rather than of search. We engage this head-on rather than against it: LeWAM's action pathway \emph{is} exactly such a predictive representation, and our own multi-seed results are consistent with the challenge, history-conditioned behavior cloning far outperforming online CEM on contact, while the fine A/B choices within the action head (detached versus action-grounded decoder, raw-action versus embedding-target, anti-collapse regularization) move success rate only within training-run noise. Our goal-conditioned head-to-head reaches the same verdict from the opposite side: an amortized, planning-free inverse-dynamics policy on the shared latent matches sampling-based search, and where both flat policies collapse together on contact-rich robomimic the limiter is what the frozen latent exposes about the task, not the choice between Markovian and history conditioning, the same representation bound MR.Q identifies. We therefore make no state-of-the-art success-rate claim. Instead we keep planning where it provably helps under the per-family rule and lean on the history-conditioned action pathway, recovering contact state, where pure search does not. The honest contribution is threefold: the per-family ``when does world-model planning help'' analysis; the past-action-conditioned action pathway that recovers contact state, where history-conditioned behavior cloning far exceeds CEM; and the reframing of a JEPA world model and its policy head as a single latent world-action model trained jointly, the latent-with-planning-analysis counterpart to a video-generative WAM such as DreamZero~\cite{dreamzero}.

% =====================================================================
% CITATION KEY -> arXiv ID + TITLE  (build the .bib from this)
% "verify" = arXiv id is after the Jan-2026 knowledge cutoff;
%            check the PDF (title/authors/numbers) before citing.
%
%   lewm            2603.19312  "Learning a Latent World Model ... (LeWM)"        [verify]
%   dinowm          2411.04983  "DINO-WM: World Models on Pre-trained Visual Features ..."
%   vjepa2          2506.09985  "V-JEPA 2 / V-JEPA 2-AC: Self-Supervised Video World Models for Action"
%   pldm            2502.14819  "PLDM: Planning with a Latent Dynamics Model (navigation)"
%   tempstraight    2603.12231  "Temporal Straightening for Latent Planning" (Wang/Bounou/Zhou/Balestriero/Rudner/LeCun/Ren) [verify, post-cutoff]
%   tdjepa          2510.00739  "TD-JEPA: Temporal-Difference JEPA for Zero-Shot RL"
%   exit            1705.08439  "Thinking Fast and Slow with Deep Learning and Tree Search (ExIt)"
%   tdmpc2          2310.16828  "TD-MPC2: Scalable, Robust World Models for Continuous Control"
%   tdmpc           2203.04955  "Temporal Difference Learning for Model Predictive Control (TD-MPC)"
%   muzero          1911.08265  "Mastering Atari, Go, Chess and Shogi by Planning ... (MuZero)"
%   pijepa          2603.25981  "PiJEPA: Policy-Initialized JEPA planning"                [verify]
%   newt            2511.19584  "Newt: Learning Massively Multitask World Models for Control" [verify]
%   mbop            2008.05556  "Model-Based Offline Planning (MBOP)"
%   gcidm           2605.08732  "Latent Geometry Beyond Search: Amortizing Planning in World Models (GC-IDM)" [verify]
%   diffuser        2205.09991  "Planning with Diffusion for Flexible Behavior Synthesis (Diffuser)"
%   decisiondiffuser 2211.15657 "Is Conditional Generative Modeling all you need ...? (Decision Diffuser)"
%   dppo            2409.00588  "Diffusion Policy Policy Optimization (DPPO)"
%   worlddp         2606.08775  "Unifying Object-Centric World Models and Diffusion Policy ... (WorldDP)" [verify]
%   ldp             2504.16925  "Latent Diffusion Planning for Imitation Learning (LDP)"
%   ffjepa          2606.09311  "FF-JEPA (feed-forward JEPA for manipulation/planning)"   [verify]
%   diffusionpolicy 2303.04137  "Diffusion Policy: Visuomotor Policy Learning via Action Diffusion"
%   vqbet           2403.03181  "VQ-BeT: Behavior Generation with Latent Actions"
%   mrq             2606.05555  "Representation Learning Enables Scalable Multitask Deep RL (MR.Q)" [VERIFIED 2026-06-20: PDF read @ proposal/papers/2606.05555.pdf]
%   openvla         2406.09246  "OpenVLA: An Open-Source Vision-Language-Action Model"
%   rt2             2307.15818  "RT-2: Vision-Language-Action Models Transfer Web Knowledge to Robotic Control"
%   dreamzero       2602.15922  "World Action Models are Zero-shot Policies (Ye et al.)"  [verify, post-cutoff]
%
% NOTE: numeric robomimic results (LDP 0.69/0.70/0.46; our histbc vs CEM=0/0/0)
% are reported in the experiments section, not here, to keep Related Work claim-free
% of our own success-rate figures (the action-head A/B is within training noise).
% =====================================================================


\section{Method: LeWAM}\label{sec:method}
\textbf{Latent and dynamics.} \todo{encoder (DINOv2/ViT) $\to z$; JEPA predictor for $z_{t+1}$ with SIGReg (isotropy $\Rightarrow$ plannable); decoder-free.}\\
\textbf{Action pathway.} \todo{autoregressive intention head on $(z_{\le t}, a_{<t})$ $\to$ action-embedding $\to$ decoder $\to$ raw action; jointly trained with the WM loss. Multitask: frozen-CLIP task token; padded+masked action dims for cross-embodiment.}\\
\textbf{Inference modes.} history-BC / intuition-guided / CEM-planning, all on one latent. \todo{equations for the loss: $\mathcal{L}=\mathcal{L}_{\text{pred}}+\lambda\,\mathcal{L}_{\text{sigreg}}+w_{\text{intent}}\mathcal{L}_{\text{intent}}+w_{\text{act}}\mathcal{L}_{\text{act}}$.}

\section{Experiments}
\textbf{Setup.} \todo{LeWM latent; robomimic Lift/Can/Square (image, PH, 200 demos) + OGBench-Cube/Push-T/Reacher/Two-Room. Eval: histbc / guided / CEM-planning, $N{=}50$, 3 seeds; per-task budgets lift100/can240/square320.}

\subsection{Planning is per-family}\label{sec:perfamily}
\begin{center}\footnotesize\setlength{\tabcolsep}{4pt}
\begin{tabular}{@{}lcccc|ccc@{}}
\toprule
& \multicolumn{4}{c|}{goal-reaching (works)} & \multicolumn{3}{c}{robomimic (contact)} \\
CEM-planning SR & Cube & Push-T & Two-Room & Reacher & Lift & Can & Square \\
\midrule
LeWAM (latent CEM) & 76 & 80 & 90 & 88 & \textbf{0} & \textbf{0} & \textbf{0} \\
\bottomrule
\end{tabular}
\end{center}
\todo{single-seed $N{=}50$ (3-seed from converging goal-reaching arms). Cube progression: $68$ (LeWM repro) $\to 78$ (planning) $\to 88$ (intuition-guided). GC-IDM~\cite{gcidm} reference on the same envs for context.}
\textbf{Mechanism.} Both cube and robomimic are grasp tasks under an identical eval, so the contrast is not ``push vs.\ grasp.'' On cube the goal-image shows the object \emph{displaced}, so reducing latent goal-cost requires grasping and carrying; on robomimic-lift the goal is the cube a few cm higher, visually swamped by arm pose, so CEM matches pose and skips the grasp. The failure has \emph{two distinct causes} across the contact family. On \textbf{Lift}, the success variable (grasp of a small cube) is only weakly exposed in a single frame---a probe decodes grasp state at AUC $0.82$ versus $0.99$ from short history (\S\ref{sec:history})---so CEM matches arm pose and skips the grasp. On \textbf{Can/Square} the single-frame latent \emph{already} decodes grasp/lift well (probe $0.96$/$0.999$), so the cause is not perception but \textbf{multi-stage structure}: a flat horizon-$H$ goal-image cost cannot decompose approach$\to$grasp$\to$transport$\to$place, the regime hierarchical planners~\cite{worlddp} target. Demonstrations are also narrow ($200$ vs $\sim$1M transitions). \todo{figure: probe AUC per task (Lift 0.82 vs Can/Square 0.96/0.999); ablate aligned-goal robomimic CEM (still $\approx 0$).}

\subsection{History-conditioned BC recovers contact state}\label{sec:history}
\begin{center}\small
\begin{tabular}{@{}lccc@{}}
\toprule
robomimic (image, PH) & Lift & Can & Square \\
\midrule
LeWAM history-BC ($N{=}50$, 3-seed) & 93.3 & 70.7 & 61.3 \\
LeWAM CEM-planning & 0 & 0 & 0 \\
LDP~\cite{ldp} (latent diffusion planner) & 69 & 70 & 46 \\
\bottomrule
\end{tabular}
\end{center}
\todo{detach=TRUE converged ckpts; the lift/can/square 3-seed shown is EVAL-seed on one trained model (train-seed robustness per \S\ref{sec:negative}); we beat LDP on Square. Mechanism is task-dependent: on \textbf{Lift} history recovers grasp state a single frame under-exposes (probe AUC $0.82$ single-frame $\to$ $0.99$ latent-history $\to$ $0.999$ full $[z,a]$), lifting post-grasp $P(\text{success}\,|\,\text{grasp})$ from $76$ to $99$; on \textbf{Can/Square} the single frame already decodes grasp ($0.96$/$0.999$), so history-BC instead supplies the sequential multi-stage execution flat CEM cannot.}

\subsection{The action-head objective is within training noise}\label{sec:negative}
\begin{center}\small
\begin{tabular}{@{}lcc@{}}
\toprule
can, $N{=}50$ & per training run & spread (all runs) \\
\midrule
detach=TRUE (embedding-intuition) & 74 / 72 / 76 $=$ \textbf{74.0} & \{70.7, 80, 74, 72, 76\} \\
solveact (raw action, $w_{\text{intent}}{=}0$) & 70 / 70 $=$ \textbf{70.0} & \{76.7, 70, 70\} \\
detach=FALSE / anti-collapse (SIGReg) & \multicolumn{2}{c}{also $\sim$72, within the same band} \\
\bottomrule
\end{tabular}
\end{center}
The variants are \textbf{statistically tied ($\sim$72$\pm$5)}: the $\sim$10-point training-run variance swamps the $4$--$7$-point effects that single runs showed. Earlier ``solveact wins'' / ``\texttt{intent\_loss} hurts'' headlines were single-training-run artifacts. We additionally verify that \texttt{act\_emb} does not collapse (its std \emph{grows} to $1.28$ without regularization) and that anti-collapse SIGReg does not help ($80$ vs.\ $76$). We therefore report this as a negative result and caution that single-training-run A/Bs (and eval-seed-only ``3-seed'') mislead.

\section{Discussion}
\todo{Limitations: contribution is analysis + mechanism, not SR. Next: hierarchical subgoal-LeWAM (WM plans subgoal latents, history policy executes) to attack the contact planning$=0$ regime (WorldDP/LDP-style); the two-part diagnosis motivates it precisely---subgoaling directly targets the \emph{multi-stage} failure on Can/Square (where the latent already sees grasp), while the \emph{perception} failure on Lift needs richer per-frame features instead. Open question: does subgoaling beat flat history-BC on the multi-stage tasks.}

\bibliographystyle{plain}
\bibliography{lewam_refs}
\end{document}
